% ID: SEL-D0051-OA-E02
% ESTADO: validada. S052.
\textbf{a)} Las tres ramas son polinomios. En $x=-3$, los límites laterales y el valor son $6$. En $x=2$, los límites laterales y el valor son $1$. Por tanto, $f$ es continua en todo $\mathbb R$. Las derivadas de las ramas son $4x$, $-1$ y $1$, respectivamente. Los cocientes incrementales laterales en $-3$ dan $-12$ y $-1$, y en $2$ dan $-1$ y $1$. En ambos puntos difieren, de modo que $f$ sólo es derivable fuera de $\{-3,2\}$.

\textbf{b)} En los dos primeros tramos la derivada es negativa. La continuidad en $-3$ une ambos descensos en uno solo. Así, $f$ decrece estrictamente en $(-\infty,2]$ y crece estrictamente en $[2,\infty)$.

\textbf{c)} El cambio de decreciente a creciente da un mínimo relativo y absoluto en $x=2$, de valor $\boxed{1}$. En $-3$ la monotonía se mantiene y no hay extremo. No hay máximo relativo.
